% !TEX program = xelatex
% =============================================================================
%  墨经 Mojing · 书籍模板示例
%  《线性代数：从方程到结构》节选（前言、第 1—2 章、答案、参考文献、索引）
%
%  编译：latexmk main.tex   （自动处理参考文献 biber 与索引 makeindex）
%  章节源文件在 chapters/ 目录中，每章一个文件。
% =============================================================================
\documentclass{mojing-book}
% 常用选项：openany（章可从偶数页起）、latin=stix（改用 Times 风格西文）、
%           accent=mjIndigo（更换主题色）、print（黑色超链接）
\addbibresource{refs.bib}

\booktitle{线性代数}
\booksubtitle{从方程到结构}
\bookentitle{LINEAR ALGEBRA: FROM EQUATIONS TO STRUCTURES}
\bookauthor{林\quad 远\quad 著}
\bookedition{排版示例 · 节选}
\publisher{墨经书局}{北京}{2026}

\begin{document}

% -----------------------------------------------------------------------------
\frontmatter
% -----------------------------------------------------------------------------
\halftitlepage
\maketitle

\begin{copyrightpage}
{\heitimd 图书在版编目（CIP）数据}\par
线性代数：从方程到结构 / 林远著. —北京：墨经书局，2026.10\par
ISBN 978-7-000-00000-0\par
Ⅰ. ①线…\quad Ⅱ. ①林…\quad Ⅲ. ①线性代数—高等学校—教材\quad Ⅳ. ①O151.2

\medskip
\noindent\rule{\linewidth}{.4pt}

\begin{tabular}{@{}ll@{}}
责任编辑 & 王\quad 某 \\
封面设计 & 赵\quad 某 \\
出版发行 & 墨经书局 \\
开\qquad 本 & $787\times1092$\quad 1/16 \\
版\qquad 次 & 2026 年 10 月第 1 版 \\
\end{tabular}

\medskip
{\color{mjGray}本页所列书号、编目数据与出版信息均为排版示例，并非真实出版物。}
\end{copyrightpage}

\unnumberedchapter{前言}

线性代数是大学数学中最早接触、也最常用到的一门课程。它起源于解线性方程组这样具体的
问题，却发展成为描述“线性结构”的普遍语言：从解析几何到微分方程，从量子力学到数据科学，
处处可以看到向量空间与线性映射的身影。

本书的写法是“从方程到结构”。我们从中国古代数学名著《九章算术》中的“方程术”讲起，
先熟练掌握消元法与矩阵运算，再逐步过渡到向量空间、线性映射、特征值等抽象概念。我们
相信，在具体计算中积累起来的直觉，是理解抽象理论最可靠的基础。

全书每章开头列出本章内容，正文中的定义、定理、例题统一按章编号，如“定理~1.4”
“例~2.7”。重要的术语在首次出现时以醒目的字体标出，并收入书末的索引。每章末尾附有
习题，分为 A、B 两组：A 组是巩固基本概念与方法的练习，B 组则需要更多的思考。书末给出了
部分习题的答案与提示。各章最后的“历史注记”简要介绍相关概念的来龙去脉，供有兴趣的读者
参阅。

本书可作为理工科各专业线性代数课程的教材，讲授约需 64 学时。

\begin{flushright}
作\quad 者\\
2026 年 10 月
\end{flushright}

\tableofcontents

\unnumberedchapter{符号说明}

\begin{center}
\renewcommand\arraystretch{1.35}
\begin{tabular}{@{}p{.24\textwidth}p{.62\textwidth}@{}}
\toprule
符号 & 含义 \\
\midrule
$\R$，$\C$            & 实数域，复数域 \\
$\R^n$                & $n$ 维实列向量的全体 \\
$\vb{x},\vb{b}$       & 列向量（黑斜体） \\
$A=(a_{ij})_{m\times n}$ & $m$ 行 $n$ 列的矩阵，$a_{ij}$ 为第 $i$ 行第 $j$ 列的元素 \\
$(A\mid\vb{b})$       & 线性方程组 $A\vb{x}=\vb{b}$ 的增广矩阵 \\
$r_i\leftrightarrow r_j$ & 交换第 $i$、$j$ 两行 \\
$cr_i$，$r_j+kr_i$    & 第 $i$ 行乘以 $c$；把第 $i$ 行的 $k$ 倍加到第 $j$ 行 \\
$A\T$                 & 矩阵 $A$ 的转置 \\
$\det A$，$\abs{A}$   & 方阵 $A$ 的行列式 \\
$S_n$，$\sgn\sigma$   & $n$ 元置换的全体，置换 $\sigma$ 的符号 \\
$I_n$                 & $n$ 阶单位矩阵 \\
\bottomrule
\end{tabular}
\end{center}

% -----------------------------------------------------------------------------
\mainmatter
% -----------------------------------------------------------------------------
\partintro{线性代数的故事从解方程开始。本部分介绍求解线性方程组的消元法，以及与之
紧密相连的行列式理论。消元法给出了判断方程组是否有解、有多少解的完整答案；行列式则把
这些答案浓缩为一个数，并揭示了它的几何意义——有向面积与有向体积。}
\part{线性方程组与行列式}

% =============================================================================
%  第 1 章  线性方程组与消元法
% =============================================================================
\chapter{线性方程组与消元法}\label{chap:systems}

\epigraph{以右行上禾遍乘中行，而以直除。}{《九章算术 · 方程》}
\chaptertoc

求解线性方程组是线性代数最基本的问题。本章介绍解线性方程组的消元法，引入矩阵和初等行
变换的语言，并完整地回答两个问题：方程组什么时候有解？有解时，解是否唯一？

\section{从《九章算术》说起}\label{sec:jiuzhang}

成书于公元一世纪前后的《九章算术》是中国古代最重要的数学著作。其第八卷名为“方程”，
专门讨论线性方程组的解法。该卷的第一题是：

\begin{quote}
\kaishu
今有上禾三秉，中禾二秉，下禾一秉，实三十九斗；上禾二秉，中禾三秉，下禾一秉，实三十四斗；
上禾一秉，中禾二秉，下禾三秉，实二十六斗。问上、中、下禾实一秉各几何？
\end{quote}

这里“禾”指带秆的谷物，“秉”意为捆，“实”是打下的谷粒。设上、中、下三等禾每秉各出谷
$x$，$y$，$z$ 斗，题意就是
\begin{equation}\label{eq:he}
\left\{
\begin{aligned}
  3x+2y+\phantom{3}z&=39,\\
  2x+3y+\phantom{3}z&=34,\\
  x+2y+3z&=26.
\end{aligned}
\right.
\end{equation}

古人用算筹在算板上布列数字：每个方程占一列，自上而下依次是各未知数的系数和常数项，
三个方程自右向左排成三列（\cref{fig:suanchou}）。这样布列出来的数表呈方形，每一列称为
一“行”，这便是“方程”一名的由来。

\begin{figure}[htbp]
\centering
\begin{tikzpicture}[x=15mm, y=8.5mm, font=\kaishu]
  \fill[mjLight] (0,0) rectangle (3,4);
  \foreach \x in {1,2} \draw[mjRule, line width=.4pt] (\x,0) -- (\x,4);
  \foreach \y in {1,2,3} \draw[mjRule, line width=.4pt] (0,\y) -- (3,\y);
  \draw[mjInk, line width=.8pt] (0,0) rectangle (3,4);
  \foreach \c/\name in {0/左行, 1/中行, 2/右行}
    \node[font=\small\heitimd, mjAccent] at (\c+.5,4.45) {\name};
  \foreach \r/\name in {3/上禾, 2/中禾, 1/下禾, 0/实}
    \node[font=\small\heitimd, mjGray, anchor=west] at (3.12,\r+.5) {\name};
  \foreach \c/\a/\b/\cc/\d in {0/一/二/三/二十六, 1/二/三/一/三十四, 2/三/二/一/三十九} {
    \node at (\c+.5,3.5) {\a};
    \node at (\c+.5,2.5) {\b};
    \node at (\c+.5,1.5) {\cc};
    \node at (\c+.5,0.5) {\d};
  }
\end{tikzpicture}
\caption{“方程”第一题在算板上的布列：每列一个方程，自右向左排列}
\label{fig:suanchou}
\end{figure}

《九章算术》给出的解法称为\term[fang cheng shu]{方程术}，其要点是\term[bian cheng zhi chu]{遍乘直除}：
用右行上禾的系数 $3$ 遍乘中行各数，再从中行连续减去右行，直到中行顶端的数变为 $0$；对左行
作同样的处理。用今天的记号，把三个方程写成数表（每个方程占一行），这一步就是
\[
\left[\begin{array}{ccc|c}
3 & 2 & 1 & 39\\
2 & 3 & 1 & 34\\
1 & 2 & 3 & 26
\end{array}\right]
\xrightarrow[\;3r_3-r_1\;]{\;3r_2-2r_1\;}
\left[\begin{array}{ccc|c}
3 & 2 & 1 & 39\\
0 & 5 & 1 & 24\\
0 & 4 & 8 & 39
\end{array}\right]
\xrightarrow{\;5r_3-4r_2\;}
\left[\begin{array}{ccc|c}
3 & 2 & 1 & 39\\
0 & 5 & 1 & 24\\
0 & 0 & 36 & 99
\end{array}\right],
\]
其中 $3r_2-2r_1$ 表示“第二行乘以 $3$ 再减去第一行的 $2$ 倍”。最后一行即 $36z=99$，于是
$z=\frac{11}{4}$；代回第二行得 $y=\frac{17}{4}$，再代回第一行得 $x=\frac{37}{4}$。也就是说，
上禾一秉九斗四分斗之一，中禾一秉四斗四分斗之一，下禾一秉二斗四分斗之三——这与原书给出的
答案完全一致。

这个两千年前的算法，与今天计算机求解线性方程组所用的方法本质上是同一个。下面我们把它
整理成一般的理论。

\section{矩阵与初等行变换}\label{sec:matrix}

\begin{definition}\label{def:system}
含 $n$ 个未知数 $x_1,\dotsc,x_n$ 的 $m$ 个一次方程组成的方程组
\begin{equation}\label{eq:general-system}
\left\{
\begin{aligned}
  a_{11}x_1+a_{12}x_2+\dotsb+a_{1n}x_n&=b_1,\\
  a_{21}x_1+a_{22}x_2+\dotsb+a_{2n}x_n&=b_2,\\
  &\;\;\vdots\\
  a_{m1}x_1+a_{m2}x_2+\dotsb+a_{mn}x_n&=b_m
\end{aligned}
\right.
\end{equation}
称为\term[xian xing fang cheng zu]{线性方程组}。使每个方程都成立的数组 $(x_1,\dotsc,x_n)$ 称为它的一个
\term[jie]{解}，全体解组成的集合称为\term[jie ji]{解集}。有解的方程组称为\term[xiang rong]{相容}的。
若 $b_1=\dotsb=b_m=0$，则称方程组是\term[qi ci xian xing fang cheng zu]{齐次}的。
\end{definition}

方程组 \cref{eq:general-system} 的全部信息都包含在系数 $a_{ij}$ 与常数项 $b_i$ 中。把它们排成矩形数表，
\[
A=\begin{bmatrix}
a_{11} & a_{12} & \cdots & a_{1n}\\
a_{21} & a_{22} & \cdots & a_{2n}\\
\vdots & \vdots &        & \vdots\\
a_{m1} & a_{m2} & \cdots & a_{mn}
\end{bmatrix},
\qquad
(A\mid\vb{b})=\left[\begin{array}{cccc|c}
a_{11} & a_{12} & \cdots & a_{1n} & b_1\\
a_{21} & a_{22} & \cdots & a_{2n} & b_2\\
\vdots & \vdots &        & \vdots & \vdots\\
a_{m1} & a_{m2} & \cdots & a_{mn} & b_m
\end{array}\right],
\]
分别称为方程组的\term[xi shu ju zhen]{系数矩阵}与\term[zeng guang ju zhen]{增广矩阵}。一般地，$m$ 行 $n$ 列的
矩形数表称为 $m\times n$ \term[ju zhen]{矩阵}。

解方程组时我们对方程所做的操作，翻译成对增广矩阵的行所做的操作，就是下面三种变换。

\begin{definition}\label{def:row-op}
对矩阵施行的下列三种变换称为\term[chu deng hang bian huan]{初等行变换}：
\begin{enumerate}
  \item \textbf{对换}：交换两行的位置，记作 $r_i\leftrightarrow r_j$；
  \item \textbf{倍乘}：以非零数 $c$ 乘某一行的所有元素，记作 $cr_i$；
  \item \textbf{倍加}：把某一行的 $k$ 倍加到另一行上，记作 $r_j+kr_i$。
\end{enumerate}
若矩阵 $B$ 可以由 $A$ 经过有限次初等行变换得到，则称 $A$ 与 $B$ \term[hang deng jia]{行等价}，记作 $A\sim B$。
\end{definition}

每种初等行变换都可以用同类型的变换复原：$r_i\leftrightarrow r_j$ 的逆是它自己，$cr_i$ 的逆是
$\frac1cr_i$，$r_j+kr_i$ 的逆是 $r_j-kr_i$。因此行等价是一个等价关系。

\begin{theorem}\label{thm:row-equivalent}
若两个线性方程组的增广矩阵行等价，则这两个方程组同解。
\end{theorem}

\begin{proof}
只需对一次初等行变换证明。对换与倍乘显然不改变解集。对倍加变换 $r_j+kr_i$：若某数组满足
第 $i$、$j$ 两个方程，则它满足“第 $j$ 个方程加上第 $i$ 个方程的 $k$ 倍”；反之，新方程组的解满足
第 $i$ 个方程和新的第 $j$ 个方程，从而满足“新的第 $j$ 个方程减去第 $i$ 个方程的 $k$ 倍”，即原来的
第 $j$ 个方程。其余方程不变，故两个方程组同解。
\end{proof}

\section{阶梯形与高斯消元法}\label{sec:gauss}

消元的目标是把增广矩阵化为一种“阶梯”形状，从而可以自下而上逐个解出未知数。

\begin{definition}\label{def:echelon}
称矩阵是\term[hang jie ti xing]{行阶梯形}的，如果
\begin{enumerate}
  \item 全为零的行（如果有的话）都在矩阵的底部；
  \item 每个非零行的第一个非零元（称为该行的\term[zhu yuan]{主元}）都在上一行主元的右侧。
\end{enumerate}
若行阶梯形矩阵还满足：(3)~每个主元都是 $1$；(4)~主元所在列的其余元素都是 $0$，则称它是
\term[jian hua hang jie ti xing]{简化行阶梯形}的。
\end{definition}

下面两个矩阵分别是行阶梯形与简化行阶梯形的典型形状，其中 $\blacksquare$ 表示主元（非零），
$*$ 表示任意数：
\[
\begin{bmatrix}
\blacksquare & * & * & * & *\\
0 & \blacksquare & * & * & *\\
0 & 0 & 0 & \blacksquare & *\\
0 & 0 & 0 & 0 & 0
\end{bmatrix},
\qquad
\begin{bmatrix}
1 & 0 & * & 0 & *\\
0 & 1 & * & 0 & *\\
0 & 0 & 0 & 1 & *\\
0 & 0 & 0 & 0 & 0
\end{bmatrix}.
\]

\begin{theorem}\label{thm:rref}
任一矩阵都与唯一的一个简化行阶梯形矩阵行等价。
\end{theorem}

存在性由下面的算法给出；唯一性的证明留作习题（第~\ref{ex:rref-unique} 题）。

\begin{algorithm}[高斯--若尔当消元法]\label{alg:gauss-jordan}
\textbf{输入}：$m\times n$ 矩阵 $A$。\quad\textbf{输出}：与 $A$ 行等价的简化行阶梯形矩阵。
\begin{enumerate}[label=\arabic*., leftmargin=2em, itemsep=1pt]
  \item 若 $A$ 为零矩阵，转第 4 步。否则找出最左边的非零列，在其中取一个非零元作为主元，
        必要时对换两行，把主元所在的行换到第一行。
  \item 用倍加变换，把主元下方的元素全部化为零。
  \item 暂时“遮住”第一行，对余下的子矩阵重复第 1、2 步，直到矩阵成为行阶梯形。
        （以上称为\emph{前向消元}。）
  \item 自下而上，对每个主元：先用倍乘把它化为 $1$，再用倍加变换把它上方的元素化为零。
        （称为\emph{回代}。）
\end{enumerate}
\end{algorithm}

\begin{example}\label{ex:free-variables}
解线性方程组
\[
\left\{
\begin{aligned}
   x_1+2x_2-\phantom{3}x_3+3x_4&=4,\\
  2x_1+4x_2+\phantom{3}x_3+3x_4&=11,\\
  -x_1-2x_2+3x_3-5x_4&=-2.
\end{aligned}
\right.
\]
\end{example}

\begin{solution}
对增广矩阵施行高斯--若尔当消元：
\begin{align*}
&\left[\begin{array}{rrrr|r}
 1 &  2 & -1 &  3 &  4\\
 2 &  4 &  1 &  3 & 11\\
-1 & -2 &  3 & -5 & -2
\end{array}\right]
\xrightarrow[\;r_3+r_1\;]{\;r_2-2r_1\;}
\left[\begin{array}{rrrr|r}
 1 & 2 & -1 &  3 & 4\\
 0 & 0 &  3 & -3 & 3\\
 0 & 0 &  2 & -2 & 2
\end{array}\right]\\
&\xrightarrow[\;r_3-2r_2\;]{\;\frac13r_2\;}
\left[\begin{array}{rrrr|r}
 1 & 2 & -1 &  3 & 4\\
 0 & 0 &  1 & -1 & 1\\
 0 & 0 &  0 &  0 & 0
\end{array}\right]
\xrightarrow{\;r_1+r_2\;}
\left[\begin{array}{rrrr|r}
 1 & 2 & 0 &  2 & 5\\
 0 & 0 & 1 & -1 & 1\\
 0 & 0 & 0 &  0 & 0
\end{array}\right].
\end{align*}
主元位于第 1、3 列，对应的 $x_1$，$x_3$ 称为\term[zhu bian liang]{主变量}，其余的 $x_2$，$x_4$ 称为
\term[zi you bian liang]{自由变量}。令 $x_2=s$，$x_4=t$，得全部解
\[
\begin{bmatrix}x_1\\x_2\\x_3\\x_4\end{bmatrix}
=\begin{bmatrix}5\\0\\1\\0\end{bmatrix}
+s\begin{bmatrix}-2\\1\\0\\0\end{bmatrix}
+t\begin{bmatrix}-2\\0\\1\\1\end{bmatrix},
\qquad s,t\in\R .
\]
\end{solution}

\section{解的存在性与唯一性}\label{sec:existence}

\cref{ex:free-variables} 中的现象具有一般性。把增广矩阵化成简化行阶梯形以后，
方程组的解的情况可以一眼看出。

\begin{theorem}\label{thm:existence}
线性方程组相容，当且仅当其增广矩阵的行阶梯形中不出现形如
$\bigl[\,0\ \ 0\ \ \cdots\ \ 0\mid b\,\bigr]$（$b\ne0$）的行，即最后一列不是主元列。
在相容的情形下：
\begin{enumerate}
  \item 若没有自由变量，则方程组有唯一解；
  \item 若有自由变量，则方程组有无穷多个解。
\end{enumerate}
\end{theorem}

\begin{proof}
由\cref{thm:row-equivalent} 与\cref{thm:rref}，不妨设增广矩阵已是简化行阶梯形。若出现
$[\,0\ \cdots\ 0\mid b\,]$（$b\ne0$），它对应方程 $0=b$，方程组无解。否则，每个非零行都可以
把主变量用自由变量表示出来；给自由变量任意赋值，主变量便被唯一确定。于是方程组有解，
且解唯一当且仅当没有自由变量。
\end{proof}

\begin{corollary}\label{cor:homogeneous}
若齐次线性方程组中未知数的个数多于方程的个数，则它必有非零解。
\end{corollary}

\begin{proof}
齐次方程组总有零解，因而相容。主元个数不超过行数 $m$，而 $m<n$，所以至少有一个自由变量，
由\cref{thm:existence}，方程组有无穷多个解，其中必有非零解。
\end{proof}

\section{几何解释}\label{sec:geometry}

两个未知数的线性方程 $a_1x+a_2y=b$（$a_1$，$a_2$ 不全为零）在平面上表示一条直线，二元线性
方程组的解就是这些直线的公共点。两条直线的位置关系只有三种：相交、平行、重合，分别对应
唯一解、无解、无穷多解（\cref{fig:lines}）。三元线性方程组的情形类似，只是直线换成了空间中的平面。

\begin{figure}[htbp]
\centering
\begin{subfigure}{.31\textwidth}
\centering
\begin{tikzpicture}[scale=.68]
  \draw[mj axis] (-.4,0) -- (3.6,0); \draw[mj axis] (0,-.4) -- (0,3.4);
  \draw[mj curve] (-.2,3.1) -- (3.3,.3);
  \draw[mj curve2] (-.2,.4) -- (3.3,3.0);
  \node[mj point] at (1.62,1.74) {};
\end{tikzpicture}
\caption{相交：唯一解}
\end{subfigure}\hfill
\begin{subfigure}{.31\textwidth}
\centering
\begin{tikzpicture}[scale=.68]
  \draw[mj axis] (-.4,0) -- (3.6,0); \draw[mj axis] (0,-.4) -- (0,3.4);
  \draw[mj curve] (-.2,2.4) -- (3.3,.0);
  \draw[mj curve2] (-.2,3.3) -- (3.3,.9);
\end{tikzpicture}
\caption{平行：无解}
\end{subfigure}\hfill
\begin{subfigure}{.31\textwidth}
\centering
\begin{tikzpicture}[scale=.68]
  \draw[mj axis] (-.4,0) -- (3.6,0); \draw[mj axis] (0,-.4) -- (0,3.4);
  \draw[mj curve, line width=2.4pt, opacity=.35] (-.2,2.9) -- (3.3,.4);
  \draw[mj curve2] (-.2,2.9) -- (3.3,.4);
\end{tikzpicture}
\caption{重合：无穷多解}
\end{subfigure}
\caption{二元线性方程组的三种情形}
\label{fig:lines}
\end{figure}

\begin{remark}
\cref{thm:existence} 说明，线性方程组的解要么不存在，要么唯一，要么有无穷多个，不会恰好有
两个或三个。这一结论也可以直接看出：若 $\vb{u}\ne\vb{v}$ 都是解，则对一切实数 $t$，
$t\vb{u}+(1-t)\vb{v}$ 也是解（习题第~\ref{ex:two-solutions} 题）。
\end{remark}

\begin{exercises}
\exgroup{A 组}
\begin{enumerate}
  \item 用消元法解下列方程组：
  \par\smallskip
  \begin{tabular}{@{}l@{\qquad}l@{\qquad}l@{}}
    (1)~$\left\{\begin{aligned}x+y+z&=6\\x-y+2z&=5\\2x+y-z&=1\end{aligned}\right.$ &
    (2)~$\left\{\begin{aligned}x_1-x_2+2x_3&=1\\2x_1-x_2+3x_3&=4\\x_1+x_3&=3\end{aligned}\right.$ &
    (3)~$\left\{\begin{aligned}x+2y&=3\\2x+4y&=5\end{aligned}\right.$
  \end{tabular}
  \item 求矩阵 $\begin{bmatrix}1&2&3\\2&4&6\\1&0&1\end{bmatrix}$ 的简化行阶梯形。
  \item 讨论 $\lambda$ 取何值时，方程组
        $\left\{\begin{aligned}x+y+\lambda z&=1\\x+\lambda y+z&=1\\\lambda x+y+z&=1\end{aligned}\right.$
        有唯一解、无解、有无穷多解，并在有解时求出全部解。
  \item （百鸡问题）《张丘建算经》载：“今有鸡翁一，直钱五；鸡母一，直钱三；鸡雏三，直钱一。
        凡百钱买鸡百只，问鸡翁、母、雏各几何？”列出方程组，并求出它的全部非负整数解。
\end{enumerate}
\exgroup{B 组}
\begin{enumerate}
  \item 证明：两个线性方程组若同解，则它们的增广矩阵不一定行等价；但若两个\emph{相容}方程组
        同解且方程个数相同，它们的增广矩阵的简化行阶梯形相同。
  \item\label{ex:rref-unique} 证明\cref{thm:rref} 中的唯一性。（提示：对矩阵的列数归纳；
        若 $R$ 与 $S$ 是 $A$ 的两个简化行阶梯形，去掉最后一列后它们仍是同一矩阵的简化行阶梯形。）
  \item\label{ex:two-solutions} 证明：若线性方程组有两个不同的解，则它有无穷多个解。
  \item 设 $A\vb{x}=\vb{b}$ 相容，$\vb{x}_0$ 是它的一个解。证明：它的全部解恰好是
        $\vb{x}_0+\vb{y}$，其中 $\vb{y}$ 取遍齐次方程组 $A\vb{y}=\vb{0}$ 的全部解。
\end{enumerate}
\end{exercises}

\begin{history}
《九章算术》的“方程术”是世界上最早的系统的线性方程组解法\parencite{guo2009}。值得注意的是，
在消元过程中，系数难免出现负数，为此“方程”章同时提出了正负数的加减法则，即“正负术”：
“同名相除，异名相益，正无入负之，负无入正之。”这也是世界上关于负数运算的最早记载之一。
公元 263 年，刘徽为《九章算术》作注，对方程术的原理作了深入的阐释，还提出了减少运算
次数的改进方法。

在欧洲，类似的消元法直到 19 世纪初才被系统使用：高斯在用最小二乘法计算小行星轨道时，
需要求解大型的线性方程组，他所用的方法今天被称为高斯消元法。算法中“回代化简”的部分以
德国测地学家 W.~若尔当（Wilhelm Jordan，1842--1899）命名，他在测量学手册中推广了这一做法；
人们常把他与法国数学家 C.~若尔当混为一谈\parencite{althoen1987}。“矩阵”（matrix）一词则由
英国数学家西尔维斯特于 1850 年引入。
\end{history}

% =============================================================================
%  第 2 章  行列式
% =============================================================================
\chapter{行列式}\label{chap:det}

\epigraph{纯数学的对象是现实世界的空间形式和数量关系。}{恩格斯《反杜林论》}
\chaptertoc

上一章用消元法回答了线性方程组的可解性问题。对于方程个数与未知数个数相同的方程组，还有
一个更紧凑的判据：方程组有唯一解，当且仅当一个由系数算出的数——行列式——不等于零。本章
从平面上的面积出发引入行列式，讨论它的性质与计算，并给出克拉默法则。

\section{面积与有向面积}\label{sec:area}

设平面上的两个向量 $\vb{u}=(a_1,a_2)$，$\vb{v}=(b_1,b_2)$ 不共线，它们张成一个平行四边形。
如\cref{fig:area}(a)，用外接矩形的面积减去四周的两对直角三角形和两个小矩形，得到
\[
  (a_1+b_1)(a_2+b_2)-a_1a_2-b_1b_2-2a_2b_1=a_1b_2-a_2b_1 .
\]
这里假定了 $\vb{v}$ 位于 $\vb{u}$ 的逆时针一侧；若交换 $\vb{u}$ 与 $\vb{v}$ 的位置，同样的式子给出
面积的相反数。于是我们定义 $\vb{u}$，$\vb{v}$ 张成的\term[you xiang mian ji]{有向面积}为
\begin{equation}\label{eq:2x2}
  D(\vb{u},\vb{v})=\begin{vmatrix}a_1&a_2\\b_1&b_2\end{vmatrix}=a_1b_2-a_2b_1 ,
\end{equation}
它的绝对值是平行四边形的面积，符号则记录了从 $\vb{u}$ 到 $\vb{v}$ 的转向。

\begin{figure}[htbp]
\centering
\begin{subfigure}{.48\textwidth}
\centering
\begin{tikzpicture}[scale=.8]
  \coordinate (O) at (0,0); \coordinate (U) at (3,1);
  \coordinate (V) at (1,2.5); \coordinate (W) at (4,3.5);
  \draw[mj grid] (0,0) grid (4,3.5);
  \fill[mj fill] (O) -- (U) -- (W) -- (V) -- cycle;
  \draw[mj thin] (0,0) rectangle (4,3.5);
  \draw[mj aux] (3,0) -- (U) -- (4,1) (1,3.5) -- (V) -- (0,2.5);
  \draw[mj curve, ->] (O) -- (U) node[midway, below right=-2pt] {$\vb{u}$};
  \draw[mj curve2, ->] (O) -- (V) node[midway, above left=-2pt] {$\vb{v}$};
  \draw[mjAccent, line width=.6pt] (U) -- (W) -- (V);
  \node[below] at (1.5,0) {\footnotesize$a_1$};
  \node[below] at (3.5,0) {\footnotesize$b_1$};
  \node[left] at (0,1.25) {\footnotesize$b_2$};
  \node[left] at (0,3) {\footnotesize$a_2$};
\end{tikzpicture}
\caption{面积 $=a_1b_2-a_2b_1$}
\end{subfigure}\hfill
\begin{subfigure}{.48\textwidth}
\centering
\begin{tikzpicture}[scale=.8]
  \coordinate (O) at (0,0); \coordinate (U) at (3,0);
  \coordinate (V) at (.6,2.2); \coordinate (V2) at (2.4,2.2);
  \fill[mj fill] (O) -- (U) -- ($(U)+(V)$) -- (V) -- cycle;
  \fill[mj fill2, opacity=.9] (O) -- (U) -- ($(U)+(V2)$) -- (V2) -- cycle;
  \draw[mj aux] (-.4,2.2) -- (5.8,2.2);
  \draw[mj curve, ->] (O) -- (U) node[midway, below] {$\vb{u}$};
  \draw[mj curve2, ->] (O) -- (V) node[midway, left] {$\vb{v}$};
  \draw[mj curve2, ->] (O) -- (V2) node[pos=.62, right=1pt] {$\vb{v}+t\vb{u}$};
  \draw[mjAccent, line width=.5pt] (U) -- ($(U)+(V)$) -- (V);
  \draw[mjSecond, line width=.5pt] (U) -- ($(U)+(V2)$) -- (V2);
\end{tikzpicture}
\caption{同底等高：$D(\vb{u},\vb{v}+t\vb{u})=D(\vb{u},\vb{v})$}
\end{subfigure}
\caption{平行四边形的有向面积}
\label{fig:area}
\end{figure}

有向面积有三条基本性质，它们都可以从 \cref{eq:2x2} 直接验证，也都有明显的几何意义：
\begin{enumerate}
  \item \textbf{多重线性}：$D$ 对每个变元都是线性的，例如
        $D(c\vb{u}+c'\vb{u}',\vb{v})=cD(\vb{u},\vb{v})+c'D(\vb{u}',\vb{v})$；
  \item \textbf{交错性}：$D(\vb{u},\vb{u})=0$，即退化的平行四边形面积为零；
  \item \textbf{规范性}：$D(\vb{e}_1,\vb{e}_2)=1$，即单位正方形的面积为 $1$。
\end{enumerate}
由 (1)(2) 立即得到 $D(\vb{u},\vb{v}+t\vb{u})=D(\vb{u},\vb{v})$，这正是\cref{fig:area}(b) 所示的
“同底等高的平行四边形面积相等”。下一节我们看到，这三条性质完全刻画了行列式。

\section{\texorpdfstring{$n$}{n} 阶行列式}\label{sec:det}

\begin{definition}\label{def:det}
设 $\det$ 是定义在全体 $n$ 阶方阵上的函数，把它看作 $n$ 个行向量的函数。若它满足
\begin{enumerate}
  \item 对每一行都是线性的（\term[duo chong xian xing]{多重线性}）；
  \item 若有两行相同，则函数值为零（\term[jiao cuo xing]{交错性}）；
  \item $\det I_n=1$，
\end{enumerate}
则称 $\det A$ 为方阵 $A$ 的\term[hang lie shi]{行列式}，也记作 $\abs{A}$。
\end{definition}

\begin{theorem}\label{thm:leibniz}
满足\cref{def:det} 的函数存在且唯一，它由下面的展开式给出：
\begin{equation}\label{eq:leibniz}
  \det A=\sum_{\sigma\in S_n}\sgn(\sigma)\,a_{1\sigma(1)}a_{2\sigma(2)}\dotsm a_{n\sigma(n)},
\end{equation}
其中 $S_n$ 是 $\{1,\dotsc,n\}$ 的全体置换；$\sigma$ 为偶置换时 $\sgn\sigma=1$，为奇置换时 $\sgn\sigma=-1$。
\end{theorem}

\begin{proof}[证明概要]
先看唯一性。由交错性与线性，交换两行使行列式变号：
$0=\det(\dotsc,\vb{a}+\vb{b},\dotsc,\vb{a}+\vb{b},\dotsc)=\det(\dotsc,\vb{a},\dotsc,\vb{b},\dotsc)+\det(\dotsc,\vb{b},\dotsc,\vb{a},\dotsc)$。
把第 $i$ 行写成 $\sum_ja_{ij}\vb{e}_j$，利用多重线性完全展开，得到 $n^n$ 项
$a_{1j_1}\dotsm a_{nj_n}\det(\vb{e}_{j_1},\dotsc,\vb{e}_{j_n})$。下标有重复的项由交错性为零；
其余各项对应置换 $\sigma$，经过若干次对换化为 $\det I_n$，得到 $\sgn\sigma$。这就证明了
函数必为 \cref{eq:leibniz}。反过来，直接验证 \cref{eq:leibniz} 满足三条公理即得存在性。
\end{proof}

当 $n=2$ 时，\cref{eq:leibniz} 就是 \cref{eq:2x2}。当 $n=3$ 时，它包含 $3!=6$ 项：
\[
\begin{vmatrix}a_{11}&a_{12}&a_{13}\\a_{21}&a_{22}&a_{23}\\a_{31}&a_{32}&a_{33}\end{vmatrix}
=a_{11}a_{22}a_{33}+a_{12}a_{23}a_{31}+a_{13}a_{21}a_{32}
-a_{13}a_{22}a_{31}-a_{11}a_{23}a_{32}-a_{12}a_{21}a_{33},
\]
可以按\cref{fig:sarrus} 所示的\term[dui jiao xian fa ze]{对角线法则}记忆。注意这一法则只适用于三阶行列式。

\begin{figure}[htbp]
\centering
\begin{tikzpicture}[x=11mm, y=8mm]
  \foreach \i in {1,2,3} \foreach \j in {1,2,3,4,5} {
    \pgfmathtruncatemacro{\jj}{mod(\j-1,3)+1}
    \ifnum\j>3
      \node[mjGray] (a\i\j) at (\j,-\i) {$a_{\i\jj}$};
    \else
      \node (a\i\j) at (\j,-\i) {$a_{\i\jj}$};
    \fi
  }
  \draw[mjInk, line width=.6pt] (.55,-.6) -- (.55,-3.4) (3.45,-.6) -- (3.45,-3.4);
  \foreach \j [evaluate=\j as \k using int(\j+2)] in {1,2,3}
    \draw[mjAccent, line width=.7pt, opacity=.75, shorten <=-3pt, shorten >=-3pt]
      (a1\j) -- (a3\k);
  \foreach \j [evaluate=\j as \k using int(\j-2)] in {3,4,5}
    \draw[mjSecond, line width=.7pt, opacity=.75, dashed, shorten <=-3pt, shorten >=-3pt]
      (a1\j) -- (a3\k);
  \node[mjAccent, font=\small] at (6.3,-1.4) {实线：取正号};
  \node[mjSecond, font=\small] at (6.3,-2.6) {虚线：取负号};
\end{tikzpicture}
\caption{三阶行列式的对角线法则：把前两列抄在右侧，沿对角线相乘}
\label{fig:sarrus}
\end{figure}

直接用 \cref{eq:leibniz} 计算高阶行列式是不现实的（$n!$ 增长极快）。实际计算依靠下面的性质，
把行列式化为三角形。

\begin{theorem}\label{thm:det-properties}
设 $A$，$B$ 是 $n$ 阶方阵，则
\begin{enumerate}
  \item $\det A\T=\det A$，因而下面关于行的结论对列同样成立；
  \item 对换两行，行列式变号；某行乘以 $c$，行列式乘以 $c$；倍加变换不改变行列式；
  \item 三角形矩阵的行列式等于主对角线上元素之积；
  \item $\det(AB)=\det A\cdot\det B$。
\end{enumerate}
\end{theorem}

性质 (2) 由定义直接得到；(1)(3)(4) 的证明见\parencite[第 2 章]{pku2013}。

\begin{example}\label{ex:det4}
计算 $D=\begin{vmatrix}1&2&3&4\\2&3&4&1\\3&4&1&2\\4&1&2&3\end{vmatrix}$。
\end{example}

\begin{solution}
各行元素之和都是 $10$。把第 2、3、4 列都加到第 1 列，再提出公因子 $10$：
\begin{align*}
D&=10\begin{vmatrix}1&2&3&4\\1&3&4&1\\1&4&1&2\\1&1&2&3\end{vmatrix}
\xlongequal[r_4-r_1]{r_2-r_1,\ r_3-r_1}
10\begin{vmatrix}1&2&3&4\\0&1&1&-3\\0&2&-2&-2\\0&-1&-1&-1\end{vmatrix}\\
&\xlongequal[r_4+r_2]{r_3-2r_2}
10\begin{vmatrix}1&2&3&4\\0&1&1&-3\\0&0&-4&4\\0&0&0&-4\end{vmatrix}
=10\times1\times1\times(-4)\times(-4)=160 .
\end{align*}
\end{solution}

\section{克拉默法则}\label{sec:cramer}

\begin{theorem}[克拉默法则]\label{thm:cramer}
设 $A$ 是 $n$ 阶方阵，$\det A\ne0$，则线性方程组 $A\vb{x}=\vb{b}$ 有唯一解
\begin{equation}\label{eq:cramer}
  x_j=\frac{\det A_j}{\det A},\qquad j=1,2,\dotsc,n,
\end{equation}
其中 $A_j$ 是把 $A$ 的第 $j$ 列换成 $\vb{b}$ 所得的矩阵。
\end{theorem}

\begin{proof}
由\cref{thm:det-properties}，初等行变换只把行列式乘以非零数，所以 $A$ 的简化行阶梯形 $R$
满足 $\det R\ne0$。$R$ 若有零行则行列式为零，故 $R=I_n$，没有自由变量，由\cref{thm:existence}
方程组有唯一解 $\vb{x}$。记 $A$ 的列为 $\vb{a}_1,\dotsc,\vb{a}_n$，则 $\vb{b}=\sum_kx_k\vb{a}_k$。
由行列式对列的多重线性与交错性，
\[
  \det A_j=\det\Bigl(\vb{a}_1,\dotsc,\sum_kx_k\vb{a}_k,\dotsc,\vb{a}_n\Bigr)
  =\sum_kx_k\det(\vb{a}_1,\dotsc,\vb{a}_k,\dotsc,\vb{a}_n)=x_j\det A ,
\]
因为 $k\ne j$ 的各项中有两列相同。
\end{proof}

\begin{example}\label{ex:he-cramer}
用克拉默法则重解\cref{sec:jiuzhang} 中的方程组 \cref{eq:he}。
\end{example}

\begin{solution}
计算四个三阶行列式：
\begin{alignat*}{2}
\det A&=\begin{vmatrix}3&2&1\\2&3&1\\1&2&3\end{vmatrix}=12, &\qquad
\det A_1&=\begin{vmatrix}39&2&1\\34&3&1\\26&2&3\end{vmatrix}=111,\\
\det A_2&=\begin{vmatrix}3&39&1\\2&34&1\\1&26&3\end{vmatrix}=51, &\qquad
\det A_3&=\begin{vmatrix}3&2&39\\2&3&34\\1&2&26\end{vmatrix}=33 .
\end{alignat*}
于是 $x=\frac{111}{12}=\frac{37}{4}$，$y=\frac{51}{12}=\frac{17}{4}$，$z=\frac{33}{12}=\frac{11}{4}$，与消元法的结果一致。
\end{solution}

\begin{remark}
克拉默法则在理论上很重要：它表明解的每个分量都是系数的有理函数。但用它来计算并不经济：
即便每个行列式都用消元法计算，也要计算 $n+1$ 个行列式，运算量约为高斯消元法的 $n$ 倍。
\end{remark}

\section{范德蒙德行列式}\label{sec:vandermonde}

\begin{theorem}\label{thm:vandermonde}
对任意数 $x_1,\dotsc,x_n$（$n\ge2$），有
\begin{equation}\label{eq:vandermonde}
  V_n=\begin{vmatrix}
  1&x_1&x_1^2&\cdots&x_1^{n-1}\\
  1&x_2&x_2^2&\cdots&x_2^{n-1}\\
  \vdots&\vdots&\vdots&&\vdots\\
  1&x_n&x_n^2&\cdots&x_n^{n-1}
  \end{vmatrix}
  =\prod_{1\le i<j\le n}(x_j-x_i).
\end{equation}
\end{theorem}

\begin{proof}
对 $n$ 归纳。$n=2$ 时显然。设结论对 $n-1$ 成立。依次把第 $n-1,n-2,\dotsc,1$ 列的 $x_1$ 倍从
第 $n,n-1,\dotsc,2$ 列中减去，第一行变为 $(1,0,\dotsc,0)$，第 $i$ 行（$i\ge2$）变为
$\bigl(1,\;x_i-x_1,\;x_i(x_i-x_1),\;\dotsc,\;x_i^{n-2}(x_i-x_1)\bigr)$。按第一行展开，再从第 $i$ 行
提出公因子 $x_i-x_1$，得
\[
  V_n=\prod_{i=2}^n(x_i-x_1)\cdot V_{n-1}(x_2,\dotsc,x_n),
\]
由归纳假设即得 \cref{eq:vandermonde}。
\end{proof}

\begin{corollary}[插值多项式的存在唯一性]\label{cor:interpolation}
设 $x_1,\dotsc,x_n$ 互不相同，则对任意 $y_1,\dotsc,y_n$，存在唯一的次数不超过 $n-1$ 的多项式
$p$，使 $p(x_i)=y_i$（$i=1,\dotsc,n$）。
\end{corollary}

\begin{proof}
设 $p(x)=c_0+c_1x+\dotsb+c_{n-1}x^{n-1}$，条件 $p(x_i)=y_i$ 是关于 $c_0,\dotsc,c_{n-1}$ 的线性方程组，
其系数矩阵的行列式正是 $V_n\ne0$。由克拉默法则，它有唯一解。
\end{proof}

例如，过四点 $(0,1)$，$(1,3)$，$(2,2)$，$(3,5)$ 的\term[cha zhi duo xiang shi]{插值多项式}是
$p(x)=1+\frac{35}{6}x-5x^2+\frac{7}{6}x^3$（\cref{fig:interp}）。

\begin{figure}[htbp]
\centering
\begin{tikzpicture}
\begin{axis}[mojing, width=.62\textwidth, height=.38\textwidth,
    xmin=-.5, xmax=3.6, ymin=-.4, ymax=6.2, xtick={1,2,3}, ytick={1,2,3,4,5},
    xlabel=$x$, ylabel=$y$, samples=100, domain=-.3:3.3]
  \addplot[mjAccent] {1+35/6*x-5*x^2+7/6*x^3};
  \addplot[only marks, mark=*, mark size=1.8pt, mjInk] coordinates {(0,1) (1,3) (2,2) (3,5)};
\end{axis}
\end{tikzpicture}
\caption{过四个点的唯一一条三次插值曲线}
\label{fig:interp}
\end{figure}

\begin{exercises}
\exgroup{A 组}
\begin{enumerate}
  \item 计算 $\begin{vmatrix}2&1&3\\1&-1&2\\3&2&1\end{vmatrix}$。
  \item 计算 $n=4$ 时的行列式 $\begin{vmatrix}a&b&b&b\\b&a&b&b\\b&b&a&b\\b&b&b&a\end{vmatrix}$。
  \item 用克拉默法则解方程组 $\left\{\begin{aligned}2x+y&=5\\x-3y&=-1\end{aligned}\right.$。
  \item 不用\cref{thm:vandermonde}，直接证明
        $\begin{vmatrix}1&1&1\\a&b&c\\a^2&b^2&c^2\end{vmatrix}=(b-a)(c-a)(c-b)$。
\end{enumerate}
\exgroup{B 组}
\begin{enumerate}
  \item 设 $A$ 是奇数阶反对称矩阵，即 $A\T=-A$。证明 $\det A=0$。
  \item 证明：平面上三点 $(x_i,y_i)$（$i=1,2,3$）共线的充要条件是
        $\begin{vmatrix}x_1&y_1&1\\x_2&y_2&1\\x_3&y_3&1\end{vmatrix}=0$。
  \item 设 $A$ 是元素都为整数的 $n$ 阶方阵。证明：对每个整数向量 $\vb{b}$，方程组 $A\vb{x}=\vb{b}$
        都有整数解，当且仅当 $\det A=\pm1$。
\end{enumerate}
\end{exercises}

\begin{history}
行列式的概念几乎同时出现在东西方。日本数学家关孝和在 1683 年的著作中已经给出了相当于
行列式展开的计算方法；十年后的 1693 年，莱布尼茨在致洛必达的信中也用到了同样的思想。
1750 年，克拉默在研究代数曲线时发表了以他命名的法则\parencite{kline1972}。范德蒙德于 1772 年
第一个把行列式作为独立的对象来研究——有趣的是，以他命名的那个行列式并没有出现在他的著作中。
1812 年，柯西系统地建立了行列式理论，证明了乘法公式 $\det(AB)=\det A\det B$，并在今天的意义下
使用了“行列式”（déterminant）这一名称；在此之前，高斯曾用这个词指二次型的判别式。
\end{history}


% -----------------------------------------------------------------------------
\backmatter
% -----------------------------------------------------------------------------
% =============================================================================
%  书末：部分习题答案与提示
% =============================================================================
\unnumberedchapter{部分习题答案与提示}

\begingroup
\small
\setlength\parindent{0pt}
\setlength\parskip{.35\baselineskip}

\subsection*{第 1 章}

{\mojingheadfont\color{mjAccent} 1.}\quad
(1)~$x=1$，$y=2$，$z=3$。\quad
(2)~第三个方程等于第二个减去第一个，有无穷多解：$x_1=3-t$，$x_2=2+t$，$x_3=t$（$t\in\R$）。\quad
(3)~消元后得 $0=-1$，无解。

{\mojingheadfont\color{mjAccent} 2.}\quad
$\begin{bmatrix}1&0&1\\0&1&1\\0&0&0\end{bmatrix}$。

{\mojingheadfont\color{mjAccent} 3.}\quad
系数行列式为 $(\lambda+2)(\lambda-1)^2$。$\lambda\ne1$ 且 $\lambda\ne-2$ 时有唯一解
$x=y=z=\frac{1}{\lambda+2}$；$\lambda=1$ 时三个方程都化为 $x+y+z=1$，有无穷多解；$\lambda=-2$ 时
三式相加得 $0=3$，无解。

{\mojingheadfont\color{mjAccent} 4.}\quad
设鸡翁、母、雏各 $x$，$y$，$z$ 只，则 $x+y+z=100$，$5x+3y+\frac{z}{3}=100$。消去 $z$ 得
$7x+4y=100$，故 $4\mid x$。令 $x=4t$，则 $y=25-7t$，$z=75+3t$，$t=0,1,2,3$，即
$(0,25,75)$，$(4,18,78)$，$(8,11,81)$，$(12,4,84)$。《张丘建算经》给出的是后三组。

{\mojingheadfont\color{mjAccent} 7.}\quad
若 $\vb{u}$，$\vb{v}$ 是解，则 $A\bigl(t\vb{u}+(1-t)\vb{v}\bigr)=t\vb{b}+(1-t)\vb{b}=\vb{b}$；
$\vb{u}\ne\vb{v}$ 时，不同的 $t$ 给出不同的解。

\subsection*{第 2 章}

{\mojingheadfont\color{mjAccent} 1.}\quad $10$。

{\mojingheadfont\color{mjAccent} 2.}\quad $(a+3b)(a-b)^3$。提示：把各列加到第一列，提出 $a+3b$。

{\mojingheadfont\color{mjAccent} 3.}\quad $\det A=-7$，$\det A_1=-14$，$\det A_2=-7$，故 $x=2$，$y=1$。

{\mojingheadfont\color{mjAccent} 5.}\quad
$\det A=\det A\T=\det(-A)=(-1)^n\det A=-\det A$。

{\mojingheadfont\color{mjAccent} 7.}\quad
若 $\det A=\pm1$，由克拉默法则，解的分量都是整数。反之，取 $\vb{b}$ 为单位向量 $\vb{e}_j$，
得整数矩阵 $X$ 使 $AX=I_n$，于是 $\det A\cdot\det X=1$，两者都是整数，故 $\det A=\pm1$。
\endgroup

\printreferences
\printindex

\end{document}
